\subsection{拓扑环的完备化}

当我们说到拓扑环的一致结构时，除非另有明确说明，我们总是指其加法群的一致结构；特别地，若 A 的加法群完备，则称 A 为完备环。

设 A 是豪斯多夫拓扑环；作为加法群，A 可以看作完备豪斯多夫交换群 \(\hat{\mathbf{A}}\) 的一个稠密子群，后者在相差一个同构的意义下被确定（§ 3，第 5 节，定理 2）。为使我们可以把 A 看作完备环的子环，必须能够把函数 \(xy\) 由连续性延拓到空间 \(\hat{\mathbf{A}} \times \hat{\mathbf{A}}\) 上。这种延拓的可能性将由下面更一般的定理得出：

定理 1. 设 E、F、G 是三个完备豪斯多夫交换群；设 A 是 E 的稠密子群，B 是 F 的稠密子群。若 f 是 A × B 到 G 中的一个连续 Z-双线性（*）映射，则 f 可由连续性延拓为 E × F 到 G 中的一个连续 Z-双线性映射。

设 \((x_{0}, y_{0})\) 是 \(E \times F\) 的任意一点，并设 \(\mathfrak{U}\)、\(\mathfrak{B}\) 分别是 \(x_{0}, y_{0}\) 的邻域滤子在 A、B 上的迹（按假设 \(\mathfrak{U}\)、\(\mathfrak{B}\) 是滤子）。为证 f 可由连续性延拓，只需证明 \(f(\mathfrak{U} \times \mathfrak{V})\) 是 G 上的一个 Cauchy 滤子基（第二章，§ 3，第 6 节，命题 11）。考虑恒等式

\[
f \left(x ^ {\prime}, y ^ {\prime}\right) - f (x, y) = f \left(x ^ {\prime} - x, y _ {1}\right) + f \left(x _ {1}, y ^ {\prime} - y\right) + f \left(x ^ {\prime} - x, y ^ {\prime} - y _ {1}\right).
\]

我们将证明，取 \((x, y)\) 与 \((x', y')\) 属于 \(\mathfrak{U} \times \mathfrak{B}\) 的一个充分小的集合，并适当选取 \(x_{1}\) 与 \(y_{1}\)，就可以使右端各项都很小。设 W 是 G 中 o 的任一邻域；由于 f 在 \((o, o) \in A \times B\) 连续，存在集合 \(U \in \mathfrak{U}\) 与集合 \(V \in \mathfrak{B}\)，使得 \(f(x' - x, y' - y) \in W\) 对 \(x \in U\)、\(x' \in U\)、\(y \in V\)、\(y' \in V\) 均成立。取一点 \(x_{1} \in U\) 与一点 \(y_{1} \in V\)；则对 U 中所有 \(x, x'\) 与 V 中所有 \(y, y'\)，我们将有

\[
f \left(x ^ {\prime} - x, y ^ {\prime} - y _ {1}\right) + f \left(x - x _ {1}, y ^ {\prime} - y\right) \in \mathrm{W} + \mathrm{W}.
\]

另一方面，偏映射 \(x \to f(x, y_1)\) 在 A 上连续；因而存在集合 \(U' \subset U\)（属于 \(\mathfrak{U}\)），使得只要 \(x\) 与 \(x'\) 属于 \(U'\)，就有 \(f(x' - x, y_1) \in W\)。同样，存在 \(V' \subset V\)（属于 \(\mathfrak{B}\)），使得只要 \(y\) 与 \(y'\) 属于 \(V'\)，就有 \(f(x_1, y' - y) \in W\)。于是，若 \((x, y)\) 与 \((x', y')\) 是 \(U' \times V'\) 的任意两点，则

\[
f \left(x ^ {\prime}, y ^ {\prime}\right) - f (x, y) \in W + W + W + W;
\]

这就证明了延拓 \(\overline{f}\)（\(f\) 的延拓）的存在性。\(\overline{f}\) 是 \(\mathbf{Z}\)-双线性的这一事实，是恒等延拓原理（第一章，§ 8，第 1 节，命题 2 的推论 1）的直接推论。

证毕。

在把这一定理应用于豪斯多夫拓扑环 A 时，我们取 E、F 与 G 为 \(\hat{A}\)，取 A 与 B 为环 A，并取 f 为 Z-双线性映射 \((x, y) \to xy\)，按假设它是连续的。我们仍用 xy 表示延拓后的函数在 \(\hat{A} \times \hat{A}\) 上的值；这个函数是 \(\hat{A}\) 上的一个合成法则，说它是 Z-双线性的，就是说它关于加法双侧分配；又由恒等延拓原理，它也是结合的。于是：

命题 6. 豪斯多夫拓扑环 A 同构于完备豪斯多夫环 \(\hat{\mathbf{A}}\) 的一个稠密子环，后者在相差一个同构的意义下被确定（并称为 A 的完备化）。

若 A 交换（相应地，有单位元），则 \(\hat{A}\) 亦然（恒等延拓原理）。

设 A 是拓扑环，不必是豪斯多夫的；设 N 是 \(\{o\}\) 在 A 中的闭包，并设 \(A' = A/N\) 是与 A 相伴的 Hausdorff 环。则 \(\hat{A}'\)——\(A'\) 的完备化——称为 A 的 Hausdorff 完备化，也记作 \(\hat{A}\)。与 §3 第 4 节命题 8 一样可以证明，A 到完备豪斯多夫拓扑环 C 中的每个连续环同态 u 都可以唯一分解为 \(u = v \circ \varphi\)，其中 v 是 \(\hat{A}\) 到 C 中的连续同态，而 \(\varphi\) 是 A 到 \(\hat{A}\) 中的典范映射。若 A、B 是两个拓扑环，且 \(u : A \to B\) 是连续同态，则存在唯一的连续同态 \(\hat{u} : \hat{A} \to \hat{B}\)，使得图表

\[
\begin{array}{c c c} \mathbf {A} & \stackrel {{u}} {{\longrightarrow}} & \mathbf {B} \\ \varphi \Big \downarrow & & \Big \downarrow \psi \\ \hat {\mathbf {A}} & \stackrel {{\hat {u}}} {{\longrightarrow}} & \hat {\mathbf {B}} \end{array}
\]

是交换的（ \(\varphi\) 与 \(\psi\) 是典范映射）；只需对 \(\psi \circ u\) 应用前面的结果。

